\section{Mean bias phenomenon}

Let $\mathcal{X} \in \mathbb{R}^{b\times s\times m}$ denote an activation tensor, where $b$ is batch size, $s$ is sequence length, and $m$ is hidden dimension. For GeMM analysis, reshape it into $\mathbf{X} \in \mathbb{R}^{l\times m}$ with $l=bs$, and write its SVD as
\[
\mathbf{X}=\sum_{i=1}^{r}\sigma_i\mathbf{u}_i\mathbf{v}_i^\top,
\qquad r=\min(l,m),
\qquad \sigma_1\ge\sigma_2\ge\cdots\ge\sigma_r>0.
\]
Define the column-wise mean vector and induced mean matrix by
\[
\mu_{\mathbf{X}}^\top=\frac{1}{l}\mathbf{1}^\top\mathbf{X}\in\mathbb{R}^{1\times m},
\qquad
\mathbf{M}_{\mathbf{X}}=\mathbf{1}\,\mu_{\mathbf{X}}^\top\in\mathbb{R}^{l\times m}.
\]
We call this phenomenon \emph{mean bias} when token representations are shifted in a common directional sense. Concretely, with normalized mean direction $\hat\mu=\mu_{\mathbf{X}}/\|\mu_{\mathbf{X}}\|_2$ and token rows $x_i$ of $\mathbf{X}$, define $p_i = x_i^\top\hat\mu.$
Across activation matrices in trained LLMs, the signs of $\{p_i\}$ are consistently concentrated on one side, indicating a coherent directional shift rather than symmetric fluctuation around zero. Figure~\ref{fig:mean_bias_definition} (left) provides one representative example from late-stage Qwen3-0.6B (layer-27 FFN input), where token-wise projection signs are nearly uniform.

We further find that the mean direction nearly recovers the leading anisotropic spike direction, which METIS identifies as the primary contributor to quantization error. Since
\[
\mu_{\mathbf{X}}=\frac{1}{l}\mathbf{X}^\top\mathbf{1}
=\frac{1}{l}(\mathbf{U}\mathbf{\Sigma}\mathbf{V}^\top)^\top\mathbf{1}
=\frac{1}{l}\mathbf{V}\mathbf{\Sigma}\mathbf{U}^\top\mathbf{1},
\]
we have the right-singular expansion
\[
\mu_{\mathbf{X}}=\mathbf{V}\alpha=\sum_{i=1}^{r}\alpha_i\mathbf{v}_i,
\qquad
\alpha_i=\frac{\sigma_i}{l}\,\mathbf{u}_i^\top\mathbf{1}.
\]
Therefore, alignment with $\mathbf{v}_1$ is large when two conditions hold jointly: (i) anisotropic spectrum with dominant $\sigma_1$, empirically observed in METIS~\citep{cao2025metis}, and (ii) limited sign cancellation in $\mathbf{u}_1^\top\mathbf{1}$ so that $|\alpha_1|$ remains large. In contrast, non-leading modes typically exhibit stronger sign alternation, making $\mathbf{u}_i^\top\mathbf{1}$ small and suppressing $\alpha_i$ for $i>1$. Consequently, $\mu_{\mathbf{X}}$ aligns much more strongly with $\mathbf{v}_1$ than with other right singular directions; empirically, this cosine similarity approaches $0.99$ (Figure~\ref{fig:mean_bias_definition}, right).

\begin{figure}[t]
\centering
\begin{subfigure}{0.49\textwidth}
\centering
\includegraphics[width=\linewidth]{figures/mean_projection_sorted.png}
\caption{Token projections on normalized mean direction.}
\end{subfigure}\hfill
\begin{subfigure}{0.49\textwidth}
\centering
\includegraphics[width=\linewidth]{figures/mean_vs_rightsv_cos_loglog.png}
\caption{Similarity between mean and right singular directions.}
\end{subfigure}
\caption{Evidence for mean-bias coherence and spectral dominance. Left: projection-sign consistency across tokens. Right: dominant alignment of mean vector with the top right singular vector.}
\label{fig:mean_bias_definition}
\end{figure}
